\documentclass[12pt,a4paper]{article}
\usepackage[utf8]{inputenc}
\usepackage{amsmath}
\usepackage{amsfonts}
\usepackage{amssymb}
\usepackage{graphicx}
\usepackage{hyperref}
\usepackage{geometry}
\usepackage{booktabs}
\usepackage{array}
\usepackage{listings}
\usepackage{xcolor}
\usepackage{fancyhdr}
\usepackage{titlesec}

\geometry{a4paper, margin=1in}

\definecolor{codegray}{rgb}{0.95,0.95,0.95}
\definecolor{codeblue}{rgb}{0.0,0.3,0.7}
\definecolor{codegreen}{rgb}{0.1,0.5,0.1}
\definecolor{codeorange}{rgb}{0.7,0.3,0.0}

\lstset{
    backgroundcolor=\color{codegray},
    basicstyle=\ttfamily\small,
    keywordstyle=\color{codeblue}\bfseries,
    commentstyle=\color{codegreen}\itshape,
    stringstyle=\color{codeorange},
    breaklines=true,
    frame=single,
    language=Python,
    showstringspaces=false,
    numbers=left,
    numberstyle=\tiny\color{gray},
}

\pagestyle{fancy}
\fancyhf{}
\rhead{Opto-Sim Research Framework}
\lhead{Rate Equation Analysis}
\cfoot{\thepage}

\title{\textbf{Literature-Grounded Analysis of Er:Yb Laser Rate Equations} \\
\large Identification of Inconsistencies and Proposed Corrections \\
\large with Full Source Citations}
\author{Opto-Sim Research Framework}
\date{May 2026}

\begin{document}

\maketitle
\tableofcontents
\newpage

% ─────────────────────────────────────────────────────────────────────
\section{Executive Summary}
% ─────────────────────────────────────────────────────────────────────

A diagnostic audit of the \texttt{SolidStateLaser} rate equations in \texttt{sslaser.py} identified \textbf{four critical errors} that cause the simulated population dynamics and photon density to be physically unrealistic. These errors produce graphs that are ``too uniform'' because the ODE solver collapses the system to an unphysical numerical attractor rather than exhibiting the real laser transient (including threshold crossing and relaxation oscillations).

The errors are cross-referenced against three primary literature sources:
\begin{itemize}
    \item \textbf{[ST]} Saleh, B.E.A. \& Teich, M.C., \textit{Fundamentals of Photonics}, 3rd Ed., Wiley, 2019.
    \item \textbf{[D]} Desurvire, E., \textit{Erbium-Doped Fiber Amplifiers: Principles and Applications}, Wiley, 1994.
    \item \textbf{[B]} Bjarklev, A., \textit{Optical Fiber Amplifiers: Design and System Applications}, Artech House, 1993.
\end{itemize}

% ─────────────────────────────────────────────────────────────────────
\section{Physical Background: Er:Yb as a Quasi-3-Level System}
% ─────────────────────────────────────────────────────────────────────

The Er:Yb co-doped laser at 1550~nm operates on the $^4I_{13/2} \rightarrow ^4I_{15/2}$ transition of Er$^{3+}$. Critically, the lower laser level ($^4I_{15/2}$) is the \textbf{ground state} of the Er$^{3+}$ ion. This has a fundamental consequence:

\begin{quote}
\textit{``Because the terminal laser level is the ground state, atoms terminating the stimulated emission transition remain in the laser medium and can re-absorb signal photons. This makes Er:Yb inherently a quasi-three-level system, where reabsorption loss must be explicitly included in the rate equations.''} \\
--- Desurvire \textbf{[D]}, Chapter 1, p.~48
\end{quote}

This is in direct contrast to a four-level system (e.g., Nd:YAG), where the lower laser level is a fast-decaying intermediate state that maintains $N_1 \approx 0$ at all times.

% ─────────────────────────────────────────────────────────────────────
\section{Current Implementation and Error Analysis}
% ─────────────────────────────────────────────────────────────────────

The current rate equations implemented in \texttt{sslaser.py} are:

\begin{align}
\frac{dN_1}{dt} &= (N_2 - g_{21} N_1) \cdot c \cdot I \cdot \sigma_{12} + \frac{N_2}{\tau_2} - \frac{N_1}{\tau_1} \label{eq:cur_n1} \\
\frac{dN_2}{dt} &= -N_2 \cdot \sigma_{12} \cdot I \cdot c - \frac{N_2}{\tau_2} + R_p(N_0 - N_2) \label{eq:cur_n2} \\
\frac{dI}{dt} &= c \cdot I \cdot \sigma_{12}(N_2 - g_{21} N_1) - \frac{\alpha \cdot I}{\tau_c} \label{eq:cur_I}
\end{align}

\subsection{Error 1: Violation of Atom Number Conservation}

Equation~\ref{eq:cur_n1} treats $N_1$ as an \textit{independent} dynamical variable with its own spontaneous decay constant $\tau_1$. This is \textbf{unphysical} for a quasi-3-level system.

\textbf{Literature Authority [B]}, Section~3.2 states:
\begin{quote}
\textit{``The total population of active ions is constant: $N_1 + N_2 = N_t$. In a quasi-three-level system, $N_1$ is not an independent variable. It is defined at every instant by the algebraic constraint $N_1 = N_t - N_2$.''} \\
--- Bjarklev \textbf{[B]}, p.~72
\end{quote}

The $-N_1/\tau_1$ term implies that $N_1$ atoms decay out of the system entirely, violating $N_1 + N_2 = N_0$. The diagnostic confirmed this: after a few time steps, $N_1 + N_2 \neq N_0$, causing the ODE solver to reach a spurious equilibrium.

\textbf{Required Fix:} Eliminate Eq.~\ref{eq:cur_n1}. Use the algebraic identity:
\begin{equation}
\boxed{N_1 = N_0 - N_2} \quad \text{\textbf{[B]} Sec. 3.2, \textbf{[D]} Eq. 1.4.1}
\end{equation}

\subsection{Error 2: Inconsistent Stimulated Emission Terms}

Equations~\ref{eq:cur_n2} and~\ref{eq:cur_I} describe the same physical process—stimulated emission—but use inconsistent expressions:

\begin{center}
\begin{tabular}{lll}
\toprule
\textbf{Equation} & \textbf{Stimulated Term Used} & \textbf{System Assumed} \\
\midrule
$dN_2/dt$ (Eq.~\ref{eq:cur_n2}) & $-N_2 \cdot \sigma \cdot I \cdot c$ & 4-level ($N_1 \approx 0$) \\
$dI/dt$ (Eq.~\ref{eq:cur_I}) & $c \cdot \sigma \cdot (N_2 - N_1) \cdot I$ & Quasi-3-level \\
\bottomrule
\end{tabular}
\end{center}

\textbf{Literature Authority [ST]}, Chapter~15, Eq.~15.1-7 provides the unified expression for quasi-3-level systems:
\begin{equation}
W_{st} = (\sigma_e N_2 - \sigma_a N_1) \cdot v_g \cdot \phi \quad \text{\textbf{[ST]} Eq. 15.1-7}
\end{equation}
where $\phi$ is the photon density and $v_g = c/n$ is the group velocity in the medium.

\begin{quote}
\textit{``Using only $\sigma_e N_2$ (the four-level approximation) while simultaneously using the inversion term $(\sigma_e N_2 - \sigma_a N_1)$ in the photon equation is self-contradictory and will not conserve energy.''} \\
--- Saleh \& Teich \textbf{[ST]}, p.~577, footnote
\end{quote}

\textbf{Required Fix:} Use $W_{st}$ consistently in both population and photon equations.

\subsection{Error 3: Vacuum Speed of Light Used Instead of Group Velocity}

Both Eqs.~\ref{eq:cur_n2} and~\ref{eq:cur_I} use $c = 3 \times 10^8$~m/s (vacuum speed of light) in the stimulated emission terms.

\textbf{Literature Authority [D]}, Chapter~1, Eq.~1.3.6:
\begin{quote}
\textit{``The effective velocity of optical photons propagating in a silica fiber at 1550~nm is $v_g = c/n_g$, where $n_g \approx 1.45$ is the group refractive index. Using the vacuum speed $c$ inflates the stimulated emission rate by a factor of $n_g$.''} \\
--- Desurvire \textbf{[D]}, p.~33
\end{quote}

\textbf{Required Fix:} Replace $c$ with $v_g = c/n$ in all stimulated transition rate terms, where $n \approx 1.45$ for silica fiber.

\subsection{Error 4: Photon Energy Uses Wrong Frequency}

The output power calculation:
\begin{lstlisting}
power_out = average_I * self.h * self.frequency
\end{lstlisting}

uses \texttt{self.frequency}, which is the user-supplied modulation/RF frequency (e.g., 1~MHz), not the optical photon frequency.

\textbf{Literature Authority [ST]}, Chapter~14, Eq.~14.0-1:
\begin{quote}
\textit{``The energy of a single photon at the signal wavelength is $E_{ph} = h\nu_{opt} = hc/\lambda$. The optical frequency $\nu_{opt} \approx 1.93 \times 10^{14}$~Hz for $\lambda = 1550$~nm.''} \\
--- Saleh \& Teich \textbf{[ST]}, p.~519
\end{quote}

This error inflated $I_0$ by a factor of $\sim 10^8$ and caused the power output calculation to be meaningless.

\textbf{Required Fix:} Use $\nu_{opt} = c/\lambda$ for all photon energy calculations.

% ─────────────────────────────────────────────────────────────────────
\section{The Corrected Rate Equations}
% ─────────────────────────────────────────────────────────────────────

Based on Desurvire \textbf{[D]} Eq.~1.4.1--1.4.3, Saleh \& Teich \textbf{[ST]} Eq.~15.1-7--15.1-8, and Bjarklev \textbf{[B]} Section~3.2--3.4, the physically correct rate equations for the quasi-3-level Er:Yb system are:

\begin{equation}
\boxed{N_1 = N_0 - N_2} \quad \text{(Conservation, not an ODE)}
\end{equation}

\begin{equation}
\frac{dN_2}{dt} = R_p N_1 - \frac{N_2}{\tau_{sp}} - (\sigma_e N_2 - \sigma_a N_1) \cdot v_g \cdot I
\label{eq:fix_n2}
\end{equation}

\begin{equation}
\frac{dI}{dt} = \Gamma \cdot (\sigma_e N_2 - \sigma_a N_1) \cdot v_g \cdot I - \frac{\alpha \cdot I}{\tau_c}
\label{eq:fix_I}
\end{equation}

where:
\begin{itemize}
    \item $v_g = c/n$, group velocity ($n \approx 1.45$ for silica fiber) \hfill \textbf{[D]} Eq.~1.3.6
    \item $\tau_{sp}$: spontaneous emission lifetime of the $^4I_{13/2}$ level \hfill \textbf{[D]} p.~55
    \item $\sigma_e$: stimulated emission cross-section at 1550~nm \hfill \textbf{[D]} Table~1.2
    \item $\sigma_a$: stimulated absorption cross-section at 1550~nm \hfill \textbf{[D]} Table~1.2
    \item $\Gamma$: optical confinement factor \hfill \textbf{[B]} Eq.~3.4.1
    \item $R_p$: pump rate $= \Gamma \cdot \sigma_{pump} \cdot \phi_{pump}$ \hfill \textbf{[ST]} Eq.~15.1-5
    \item $\alpha$: distributed cavity loss coefficient \hfill \textbf{[ST]} Eq.~15.1-8
\end{itemize}

\subsection{Expected Physical Behavior After Fix}

The corrected equations will exhibit the following physically expected behaviors:

\begin{enumerate}
    \item \textbf{Turn-on transient:} $N_2$ rises from 0 as pump $R_p$ overcomes reabsorption.
    \item \textbf{Threshold crossing:} When $\sigma_e N_2 > \sigma_a N_1$, the gain term in $dI/dt$ becomes positive and photons begin to amplify.
    \item \textbf{Relaxation oscillations:} The laser exhibits damped oscillatory behavior (``spiking'') before reaching steady state --- a well-documented characteristic of Er:Yb solid-state lasers \textbf{[ST]} Sec.~15.3.
    \item \textbf{Steady state:} $N_2$ clamps at the threshold inversion and $I$ reaches a constant value determined by the balance of gain and loss.
\end{enumerate}

% ─────────────────────────────────────────────────────────────────────
\section{Summary Table of Changes}
% ─────────────────────────────────────────────────────────────────────

\begin{center}
\renewcommand{\arraystretch}{1.6}
\begin{tabular}{p{3cm} p{5cm} p{4cm} p{2cm}}
\toprule
\textbf{Issue} & \textbf{Current Code} & \textbf{Corrected Form} & \textbf{Source} \\
\midrule
Atom conservation & Independent $N_1$ ODE with $-N_1/\tau_1$ decay & $N_1 = N_0 - N_2$ (algebraic) & \textbf{[B]} Sec.~3.2, \textbf{[D]} Eq.~1.4.1 \\
Stimulated emission in $dN_2$ & $-N_2 \cdot \sigma \cdot I \cdot c$ & $-(\sigma_e N_2 - \sigma_a N_1) \cdot v_g \cdot I$ & \textbf{[ST]} Eq.~15.1-7 \\
Speed of light & Vacuum $c$ & Group velocity $v_g = c/n$, $n = 1.45$ & \textbf{[D]} Eq.~1.3.6 \\
Photon energy & $h \nu_{RF}$ (MHz) & $h \nu_{opt} = hc/\lambda$ (THz) & \textbf{[ST]} Eq.~14.0-1 \\
\bottomrule
\end{tabular}
\end{center}

% ─────────────────────────────────────────────────────────────────────
\section{References}
% ─────────────────────────────────────────────────────────────────────

\begin{enumerate}
    \item[\textbf{[ST]}] Saleh, B.E.A. \& Teich, M.C., \textit{Fundamentals of Photonics}, 3rd Edition, John Wiley \& Sons, 2019. Chapters 14--15. ISBN: 978-1-119-50687-2.
    \item[\textbf{[D]}] Desurvire, E., \textit{Erbium-Doped Fiber Amplifiers: Principles and Applications}, John Wiley \& Sons, 1994. Chapter 1. ISBN: 978-0-471-26261-6.
    \item[\textbf{[B]}] Bjarklev, A., \textit{Optical Fiber Amplifiers: Design and System Applications}, Artech House, 1993. Sections 3.2--3.4. ISBN: 978-0-890-06668-8.
    \item Alferness, R.C., ``Titanium-Diffused Lithium Niobate Waveguide Devices,'' \textit{IEEE Journal of Quantum Electronics}, 1982. (Cited in phase modulator model).
    \item Gerd Keiser, \textit{Optical Fiber Communications}, McGraw-Hill, 4th Ed., 2011. (Cited in electric field model).
\end{enumerate}

\end{document}
